\section{Random Walk}

\begin{definition}[Random/Stochastic Process]
    \label{def:random-stochastic-process}%
    A \emph{random/stochastic process} is a sequence of random variables \(\para{X_n}_{n \in \NN}\).
\end{definition}

\begin{definition}[Random Walk]
    \label{def:random-walk}%
    A \emph{random walk} is a random process that can be expressed in the following way:
    \[
        X_n = x + \sum_{i = 1}^{n} Y_i,
    \]
    where \(x\) is a deterministic number, and \(\para{Y_i}_{n \in \NN}\)s are independent and identically distributed (i.i.d.). The random variables \(Y_i\)s the \emph{steps} of the random walk.
\end{definition}

\begin{definition}[Simple Random Walk]
    \label{def:simple-random-walk}%
    A random walk is \emph{simple} if the steps can only take values \(\pm 1\).
\end{definition}

In this case, we may set \(\Prob\para{Y_i = 1} = p\) and \(\Prob\para{Y_i = -1} = q\), with \(p + q = 1\) and \(p, q \geq 0\).

\begin{notation}
    We denote \(\Prob_x\para{A} = \Prob\para{\Conditional{A}{X_0 = x}}\) and \(\Expt_x\brac{A} = \Expt\brac{\Conditional{A}{X_0 = x}}\).
\end{notation}

\begin{example}[Gambler's Ruin Estimates]
    \label{ex:gambler-ruin-estimates}%
    We may think of \(\para{X_n}\) as the fortune of a gambler, who starts with \(x\) in their pocket at time \(0\), and every time step they bet \(1\), which he loses with probability \(\para{1 - p}\), and gets it back doubled with probability \(p\). The game ends if they reach \(0\) (which means they go bankrupt) or they reach \(a\), whichever happens first.

    We want to know the probability they reach \(a\) before going bankrupt. We let
    \[
        h\para{x} = \Prob_x \para{\para{X_n} \text{ reaches }a\text{ before reaching }0}.
    \]

    Note that \(h\para{0} = 0\) and \(h\para{a} = 1\). We have
    \begin{align*}
        h\para{x} &= \Prob_x \para{X \text{ reaches }a\text{ before }0} \\
        &= \Prob_x \para{\Conditional{X \text{ reaches }a\text{ before }0}{Y_1 = 1}} \cdot p + \Prob_x \para{\Conditional{X \text{ reaches }a\text{ before }0}{Y_1 = -1}} \cdot \para{1 - p} \\
        &= p \cdot h \para{x + 1} + \para{1 - p} \cdot h \para{x - 1}.
    \end{align*}

    We now consider particular cases of \(p\) and \(q\).

    \begin{itemize}
        \item If \(p = q = \frac{1}{2}\), this is called \emph{simple symmetric random walk}, and the system reduces to \(h\para{x} = \frac{1}{2} h \para{x + 1} + \frac{1}{2} h \para{x - 1}\).
        
        Notice that simplifying this gives
        \[
            h\para{x + 1} - h\para{x} = h\para{x} - h\para{x - 1} = \cdots = c.
        \]

        Therefore, we have \(h\para{x} = \sum_{i = 1}^{x} \para{h\para{i} - h\para{i - 1}} = cx\). Since \(h\para{a} = 1\), we have \(c = \frac{1}{a}\), so \(h\para{x} = \frac{x}{a}\).

        \item If \(p \neq q\), we first try a solution of the form \(\lambda^x\) for some \(\lambda\). We have
        \[
            \lambda^x = p \cdot \lambda^{x + 1} + q \cdot \lambda^{x - 1}
        \]
        and hence \(\lambda = 1\) or \(\lambda = \frac{q}{p}\). Therefore, the general solution gives
        \[
            h\para{x} = A + B \cdot \para{\frac{q}{p}}^x.
        \]

        Putting in the boundary conditions, we have
        \[
            h\para{x} = \frac{\para{\frac{q}{p}}^x - 1}{\para{\frac{q}{p}}^a - 1}.
        \]
    \end{itemize}
\end{example}

\begin{example}[Time to Absorption]
    \label{ex:time-to-absorption}%
    We use a similar setup as \autoref{ex:gambler-ruin-estimates}. Let the \emph{time to absorption} be
    \[
        T = \min \set{n \geq 0}{X_n \in \brce{0, a}}.
    \]

    \(T\) is the first \(n\) such that \(X_n\) is either \(0\) or \(a\). In other words, \(T\) is the duration of the game.

    We aim to find \(\Expt\brac{\Conditional{T}{X_0 = x}} = \Expt_x\brac{T}\). Let this be \(k\para{x}\).

    Using the law of total expectation, we have
    \begin{align*}
        k\para{x} &= \Expt_x\brac{\Conditional{T}{Y_1 = 1}} \cdot \Prob\para{Y_1 = 1} + \Expt_x\brac{\Conditional{T}{Y_1 = -1}} \cdot \Prob\para{Y_1 = -1} \\
        &= p \cdot \para{1 + k\para{x + 1}} + \para{1 - p} \cdot \para{1 + k\para{x - 1}} \\
        &= 1 + p \cdot k\para{x + 1} + q \cdot k\para{x - 1}.
    \end{align*}

    The initial conditions are \(k \para{0} = k \para{a} = 0\).
    \begin{itemize}
        \item If \(p = q = \frac{1}{2}\), we experiment with a solution of the form \(A x^2\). We have
        \[
            A x^2 = 1 + p A\para{x + 1}^2 + q A\para{x - 1}^2.
        \]

        Hence, \(A = -1\), and the general solution is
        \[
            k\para{x}= -x^2 + Bx + C.
        \]

        With the initial conditions, we have
        \[
            k\para{x} = x \para{a - x}.
        \]

        \item If \(p \neq q\), we experiment with a solution of the form \(Cx\). We have
        \[
            C = \frac{1}{q - p}
        \]
        giving general solution
        \[
            k\para{x} = A + \frac{1}{q - p} \cdot x + B \cdot \para{\frac{q}{p}}^x,
        \]

        Hence, putting in the initial conditions, we have
        \[
            k\para{x} = \frac{1}{q - p} \cdot x - \frac{a}{q - p} \cdot \frac{\para{\frac{q}{p}}^x - 1}{\para{\frac{q}{p}}^a - 1}.
        \]
    \end{itemize}
\end{example}